\typesetchapter{Chapter 2}{Mirror}

\noindent You will want to know how a boy raised to worship stone became a man who could take it apart. It did not happen in the Tomb, Paul, whatever your notes say. It happened the morning after they cut my hand, in a room full of mirrors, over a bowl of water no one thought to guard. Nobody hands you the thing that ruins you. They teach you a smaller thing, carefully, and look away at the wrong moment.\par

\ornament

I came to the eastern practice halls before first work with an empty stomach, a bandaged thumb, and the private conviction that I would not be made to look foolish twice.\par

The day before, I could have passed those thresholds only as Kailan, a boy under instruction, tolerated near the outer rooms because his teachers had vouched for him. That morning the first eastern gate did not ask for tolerance. It asked for record. An attendant in pale grey read my new name from a tablet where the letters had been cut so recently that dust still gathered inside them, and my name in a stranger's mouth sounded less like honour than like a key that had already turned in locks I had not seen.\par

The stone changed under my sandals as I crossed from common passage into Light. It became smoother, colder, cut at shallow opposing angles, so that a step did not merely sound but returned a faint image from below. Even the darkness was tended here, assigned rather than left to gather where it pleased. A corridor beyond the gate was crossed by thin bars of light that fell from slits too high to see and vanished into black cuts on the far wall. Dust moved through them, visible, invisible, visible again, each mote existing only where the beam permitted it.\par

I slowed despite myself. Faith reached its chambers by obedience, Balance by measure, Shield by permission, Justice by summons. Light, I had always believed, worked differently. Light showed. A man did not need to believe a lamp; he had only to open his eyes.\par

Then Va-Raedin's voice from the Hall returned to me. \textit{Do not admire the flame. Ask where the air enters.} I looked again at the shafts crossing the passage. The obvious thing was their precision. The better question was where they began, who had narrowed them, why half of them were received by slots and not stone, and what waited in the darkness that had been denied the beam. I had not answered it by the time I reached the hall.\par

The Hall of Light was not one hall, though it had one name. From the north it seemed broad and shallow; from the eastern stair it narrowed to a blade of space, mirrors along one side, black apertures cut into the other. From the candidate entrance, where I stood, it looked almost empty. Then it began to give itself away. Mirrors stood everywhere, but none where a child would have set them, tall strips no wider than a hand, small round discs at knee or shoulder, and a few that were not mirrors at all but dark polished stone that returned shape without obedience. Arcs and crosses and incomplete circles had been cut into the floor. Three doorways stood at the far end. One open and dark, one closed by bronze, one half-open with a line of brightness along its edge that I distrusted at once.\par

Va-Raedin stood at the centre, hands folded behind his back, pale grey robe, no ornament but the narrow silver thread at the cuff. He looked rested. I resented it. The ceremony had lasted half a day, and he had stood through it without shifting, leaning, or submitting to any of the small disgraces by which bodies confess themselves.\par

``You are early,'' he said.\par

``You said before first work.''\par

``I did.'' He looked at my bandage. ``Does the thumb hurt?''\par

``No.''\par

``It should.''\par

``It does a little.''\par

``Then begin with truth where truth is inexpensive.''\par

Heat moved into my face, and I said nothing, which I would learn was the only reply that did not cost me more.\par

He turned to the three doors. ``Which door opens?''\par

I nearly answered quickly, and disliked myself for the impulse. The right doorway's brightness was so plainly a lure that the lure might be to distrust it too easily. I studied the floor instead. A thin strip of light lay at the base of the left doorway; the right's brightness touched nothing; the bronze was too solid, too central, too theatrical.\par

``The left,'' I said. ``The light at its base belongs to an opening. The right only reflects one. The middle is meant to make the eye feel opposed.''\par

``Open it.''\par

I crossed the hall and pushed. Nothing. Pulled. Nothing. I searched for latch, seam, hinge, weight. Smooth stone answered me everywhere, and behind me Va-Raedin did not speak, which I hated more than correction. I pressed the bronze door once, because anger sometimes insists on humiliating itself with the body, and it gave me nothing either.\par

``You asked which door opens,'' I said. ``Not which door appears open.''\par

``No.''\par

I looked again. Doors, mirrors, floor. The nearest cut circle was worn on one side only, as though candidates had learned to step into it from the same direction. I stood inside it, and the room changed, not by moving but by being seen from the point for which it had been arranged. The mirrors on the left no longer held the three doors. They caught a plain stone panel behind me, flush with the wall, unmarked by light. I found the seam by touch and pressed, and it opened inward.\par

I tried not to smile. He saw me trying.\par

``You are pleased because you found the fourth door.''\par

``I found the door that opens.''\par

``You found it after failing.''\par

``I corrected the failure.''\par

``That is better than defending it. It is not the same as succeeding.'' He came beside me. ``What was the lesson?''\par

``That light deceives.''\par

``No. That is what boys say when a lamp embarrasses them. The lesson is that sight depends on position. From the entrance you saw three doors and chose among them. From the circle you saw the room that had arranged the choice. The hidden door was not hidden because stone wished to be secret. It was hidden because you accepted the first place given to your eyes.''\par

I have thought about that sentence in more rooms than I can count, Paul. In yours, too. Every interrogation is a man deciding which place to give your eyes.\par

``A Candidate learns where light falls,'' Va-Raedin said. ``A Candidate who follows brightness serves whoever placed the lamp.''\par

``A Veil learns who placed it,'' I said.\par

``Eventually. If he stops being proud of finding doors built to be found.''\par

The correction landed because it was deserved, and I went into the dark chamber before I could be told, partly to escape my own face.\par

\ornament

The second room was narrower and colder, lit through slits so thin they should not have lit anything, so that the chamber glimmered in fragments. On a stone table lay a mirror shard, a black measuring rod, a red glass bead fixed to a white tile, a bowl of water, two bronze wedges, and a square of cloudy glass. Va-Raedin closed the hidden door, and the darkness thickened, then organised itself around the slits.\par

``Move the red point onto the far wall.''\par

The far wall was blank.\par

``May I touch the lamp?''\par

``No.''\par

``The table?''\par

``You may touch what is on the table.''\par

I began with the shard. The nearest slit sent a fine blade of light too high to catch, and I angled the shard until the light jumped downward, bright and disobedient, crossed the black rod, and vanished. I moved the rod, then understood it was not an obstacle but a measure, its shadow reading the beam's angle more precisely than my eye. A bronze wedge held the mirror steady. The light struck the cloudy glass and widened, woke the red bead into a smeared glow, and I sharpened it until it slid up the wall and vanished again.\par

For minutes I worked in silence, at an exercise simple to understand and humiliating to perform. Each time the point steadied, the smallest movement of my wrist sent it running. My bandaged thumb made the shard hard to grip; I compensated with the wrong fingers, overcorrected, lost the beam, found it, and heard Va-Raedin breathe once through his nose. Not a sigh. Worse. At last I set the mirror flat and used the water, tipping the bowl until its surface caught the beam and threw it upward, unstable but bright, and the red point appeared on the wall, alive to every tremor in the bowl. I held my breath until the surface stilled.\par

``Acceptable,'' Va-Raedin said. I felt the word as insult and reward at once. ``Why was the water there?''\par

``To reflect.''\par

``The mirror reflects.''\par

``To reflect differently.''\par

``Everything reflects differently. Be precise.''\par

I looked into the bowl. The surface still carried the small ripples of my hand, and the red point trembled in sympathy. ``The mirror keeps its angle,'' I said. ``Water answers after the hand has stopped touching it.''\par

He was silent long enough for the answer to become visible to me as well.\par

``Better.'' That word struck harder than the first. ``Water remembers disturbance. Not as Faith uses the word. Not grief, not mercy, not witness. Motion. A careless hand leaves a pattern after the hand is gone, and light striking the pattern tells you something happened.''\par

``So what is unseen can be known by what it disturbs,'' I said.\par

``Sometimes. The water tells you disturbance occurred. It does not always tell you what disturbed it. Young minds love to turn traces into stories.''\par

I thought of Ka-Elun walking the Labyrinth without witness. Faith said God had guided her. Light would ask what marks were left by feet, breath, oil, scraped stone, fear, and time. Neither answer cancelled the other, and that unsettled me more than either would have alone.\par

He removed the shard from the table. ``Again. Without this.''\par

``That was the main tool.''\par

``It was the easiest tool.''\par

I failed. Twice I spilled water and was made to wipe the table with my sleeve, because, he said, Balance would not be summoned to clean Light's impatience. Once I lost the point entirely. Once I threw a wide inverted stain across the wall, more wound than mark. Frustration shortened my breath and made my fingers too strong, and he corrected as if time were one more instrument in the room. He set two fingers to the inside of my wrist --- cold, exact, the way one steadies an instrument that has begun to lie --- and held it there until the tremor I had been pretending not to feel stood plain against them. ``You shake,'' he said, and made it a finding, not a mercy. ``Too much angle. Too much hand. You are moving the glass because you dislike moving yourself. Stop chasing the image. Trace the path. Source, aperture, surface, eye. Again.'' By the seventh attempt I stopped forcing the point and moved around the table until slit, water, glass, tile, and wall made a relation I could feel before I could state it. The point appeared low on the wall, smaller than before, and stable. I did not step back. I waited.\par

``Acceptable.''\par

The word stung less.\par

\ornament

He has been describing the bowl of water --- the beam thrown up at the wall, the ripples answering his hand --- and I let the account run until it brushes a thing I know better than any servant of record should. Then I do the one thing this room does not forgive. I help him.\par

``The water did more than throw the beam back,'' I say. ``When the light went down into it and rose again to the wall, it did not rise where it entered. It bent at the skin of the water, and bent further the deeper it reached. The point ran from the boy because he was chasing the bend, not only his hand.'' I hear my own voice fill the space the way a voice fills a room it has measured, and I stop it a breath too late.\par

He does not move. But something goes still in him, and I watch him set the moment down somewhere behind his eyes --- exactly, in its measure, the way he sets his own name down so he can find it again. A servant keeps the account. A servant does not know, unasked, which way light bends beneath water.\par

``You have done this,'' he says. He is too careful to make it a question.\par

``I kept records once for men who ground glass,'' I say. It is a lie, and a clean one, and it does the work a lie is for --- it gives the strangeness a smaller shape to sit in. A man who counted another's lenses in a ledger is a lesser thing than a man who bends light for his living, and he lets the lesser thing be true, because he has no reason yet to build the larger one.\par

``Show me a ground glass,'' he says.\par

There is a reading lens on the side table, thick at the centre, thinned to nothing at the rim, in a handle of horn. I pass it to him. He holds it over the machine and the two reels swell, slow and enormous; he lifts it until they swim and break, and a thing moves in his face that I do not name for him --- the boy at the basin, I would guess, meeting in ground glass the thing he found in the dark and believed was his alone.\par

``It gathers light to a point,'' I say. True, and true is the cheap answer here, so I give it whole. ``The eye is one also --- a wet stone that bends the world to a point at the back of the skull.''\par

``A made eye. Held in the hand.'' He returns it with a care I do not think he knows he shows. ``We were given only the found ones. Water, oil, the flaw in a jar --- a bend a man could swear had not been arranged.'' He turns the reels back with one finger. ``Go on,'' he says, before I can. ``There were other rooms.''\par

\ornament

We passed through three more chambers before first work began elsewhere in the Mountain, and I heard the day waking beyond Light, carts toward the food galleries, voices near the washing pools, the metal rhythm of a gate being tested by Shield. In the eastern halls morning did not arrive; it was admitted. A corridor I had thought dark became pale; a pale one became useless.\par

In one chamber he made me cross a floor divided into light and shadow without stepping on lit stone, and when I finished quickly he changed nothing except my starting place, and from the new angle every certainty failed.\par

``Shadow is not darkness,'' he said.\par

``An absence of light.''\par

``Too broad.''\par

``An interruption.''\par

``Better. By what?''\par

A lamp, a slit, a pillar, my body, his body, the raised edge of a tile: each interruption made a different shape. ``By whatever enters the path,'' I said.\par

``Of light, of view, of expectation. If a man stands between you and a door, you call the door hidden. The door has not changed. The man has entered your seeing.''\par

``That seems obvious.''\par

``Most useful things do, until someone uses them against you.''\par

He crossed to a blackened panel and laid his hand flat upon it. ``Where is the exit?'' The entrance behind us was one answer and therefore probably not the answer; his hand drew my eye to the panel, too large, too flat, no seam, unless soot hid the seam. I walked toward it. He removed his hand, and the panel became simply black stone, and I stopped.\par

``Why did you come here?''\par

``You drew attention to it.''\par

``Yes.''\par

``You made yourself the lamp.''\par

``No. I made myself the obstruction.''\par

I looked back. When he had stood at the panel, his body had blocked a low mirror behind him. With him gone, the reflection crossed the room and struck the left arch, showing that the floor beyond it did not continue. The false corridor had declared itself only when he stepped away.\par

``You hid the falsehood by standing in front of the light,'' I said.\par

``I prevented you from seeing the condition that made the falsehood visible. That is not the same thing.''\par

``The distinctions are convenient.''\par

``They are survival.''\par

``To whom?'' The question came sharper than I intended.\par

He looked at me. ``To anyone who does not wish to be governed by the nearest brightness.''\par

I wanted to say Light should stand apart from the others; that where Faith clothed uncertainty in prayer and Balance in measure and Shield in necessity and Justice in record, Light at least should reveal things as they were.\par

``Light is not truth,'' he said, before I had shaped it into a question. ``It is a condition under which truth may become visible, falsehood persuasive, or danger manageable. A lamp in the wrong place blinds more thoroughly than darkness.''\par

``Then what is our charge?''\par

``To know the difference before others suffer for it.''\par

I did not ask how anyone could know the difference if Light itself could be arranged into lies. He would have approved the silence, which made keeping it mildly intolerable.\par

\ornament

There was a corridor lined with dull metal plates, each angled a little differently, that bent twice while reflections made it appear straight, and when I looked back I could no longer find the chamber we had left. I stopped. He went five steps on and paused.\par

``Do not trust backward sight merely because it comforts you.''\par

``It shows the hall.''\par

``Which hall?''\par

I studied the reflection, and the longer I looked the more wrongness surfaced. The third doorway reversed. The eastern mirror strips too high. The cut circle gone. It was not the first hall but a likeness of it, assembled from enough correct pieces to soothe the eye.\par

``Who would build a false view behind a candidate?'' I asked.\par

``Someone who knows candidates look back when uncertain.''\par

The corridor opened into a chamber lit by a single lamp on the floor, ringed by six mirrors on slender bronze stems, each facing the lamp at a different angle. Cut into the far wall, in plain letters:\par

\begin{scripture}
WHERE LIGHT FALLS IS ONLY THE FIRST QUESTION.
\end{scripture}

He adjusted one mirror by less than the width of a fingernail. Three patches of light shifted; one vanished; another woke where the wall had been dark. ``Place them so the lamp finds the door.''\par

``There is no door.''\par

``No visible door.''\par

I wasted minutes on the brightest patch before I understood it was bait, a reflection of a reflection, brightness doubled and direction corrupted; when I moved its mirror the whole arrangement collapsed into near-dark, and I reset it badly, and he did not help. I tried again, tracing the order this time, lamp to mirror, mirror to wall, wall to absence, absence to edge. My pride objected to crouching before my sense accepted that height was another form of ignorance. From the floor I saw that one mirror reflected not the lamp but the gap between two brighter paths. I shifted it toward the gap. A pale line sharpened on the wall; another mirror sent it downward; a seam appeared at ankle height. Not a door for a man. A panel. I stood too quickly, struck my shoulder against a bronze stem, and the mirror rang, and three lights jumped and vanished.\par

He looked at the trembling mirrors. I closed my eyes.\par

``You may begin again.''\par

The worst punishment was a room that merely showed you what haste had done. By the time the panel opened, my knees ached and the linen on my thumb had darkened with new blood. Inside lay a small stone token marked with the sign of Light: not a lamp, as children draw it, but a narrow aperture cut through a circle.\par

``Take it,'' he said. It was heavier than it looked. ``Proof that you opened the room.''\par

``For whom?''\par

``For you, if you are vain. For me, if you are honest.''\par

``I opened it.''\par

``You did.''\par

``You expected me to fail.''\par

``I expected you to fail several times. You did not disappoint me.''\par

I smiled then, despite myself, and he saw it and did not punish it.\par

\ornament

A passage overlooked a lower work hall through narrow slits, and I saw that nothing there belonged to Light alone. Oil came by Balance, sealed and checked. Doors by Shield. Witness by Justice. Below, Faith singers rehearsed near a side arch, and no full song reached me through the stone, only a low pressure in the rail beneath my hand, as if the wall had taken the note and returned its body without its words. A thin mist rose behind the singers. It did not come from the floor. It emerged from a line of holes above the arch, so fine I had first taken them for pores in the rock, and a hidden lamp turned the vapour pale at the exact moment the singers lifted their faces.\par

Va-Raedin let me look. After a while: ``What do you see?''\par

``A rehearsal.''\par

``What else?''\par

``Water in the air. Light placed to make it visible. Timing.'' Below, a singer raised both hands, and the mist thickened with the movement as though the prayer had called it. The stone under my palm answered a breath after the bodies below. ``Sound moving where sight cannot follow,'' I said.\par

He nodded once. And a thing entered me there, small and sharp, that I did not have the courage to say aloud: that a prayer could be sincere and placed with precision at the same time. That a blessing could be engineered without becoming false. That water could be counted in one chamber and appear as mercy in another. If the people were comforted, was the comfort a lie? If God provided through valves, did the valve disprove Him, or serve Him? The mist thinned. The stone released its low answering hum. It had been beautifully done, and I hated how beautiful it was.\par

\ornament

We came back to the upper hall by a route that did not feel like return, and the first chamber had changed in our absence, or I had. The three false doors seemed less convincing. The mirrors seemed less like objects than decisions.\par

``Again,'' he said. ``The first test.''\par

``I know the answer.''\par

``That will make the failure more instructive.''\par

He set me at the original entrance, and the room appeared as before, three doorways, the hidden panel invisible from that place. I did not go to the doors. I found the worn circle, stepped in, turned, pressed the wall. Nothing. The seam was gone. I kept my face still with effort and looked again, and saw that the circle's wear was false, darkened with rubbed dust, and the true wear lay two steps behind it, faint under a reflected shadow. I moved to it. The room shifted; the hidden panel was no longer behind me but to my left, disguised by darkness thrown from a mirror. I touched it. It opened. This time I smiled first.\par

``Better,'' he said. Not \textit{acceptable}. \textit{Better}. I stored the word away like contraband.\par

\ornament

Near the end he led me to a narrow balcony above one of the external light shafts. Cold air rose through it, dry and mineral, with a far smell of desert stone beneath. A beam fell through it, pale and almost solid, until it struck a suspended plate of polished metal, and from that plate it divided through more apertures than I could count, carved throats waiting to accept and divide what little Heaven gave.\par

``When you were a child, what did Faith tell you about the Mountain's light?''\par

``That God remembered us beneath stone.''\par

``Did you believe it?''\par

``Yes.''\par

``And now?''\par

I looked at the beam, the plate, the apertures. ``I think someone had to build this.''\par

``That is not an answer.''\par

``No,'' I said. ``But it is the question I have.''\par

He turned his head slightly. Approval, or warning; with Va-Raedin the two often wore the same face. ``Good. But understand the danger. Most children stop at miracle. Most fools stop at mechanism. Light permits neither laziness. A miracle may have a mechanism. A mechanism may still serve a miracle. What matters is who knows, who does not, and what each ignorance protects.''\par

I watched the beam strike the plate. Around the edge of the metal the light misbehaved. I leaned in, and for a moment thin bands appeared, pale, dark, pale, dark, too delicate to name before they vanished when I shifted my head.\par

``What was that?''\par

``What?''\par

``The edge. The light changed.''\par

``Edges do that.''\par

``How?''\par

He looked at the plate, then at me. ``You are not ready for every how.''\par

It would be years before I learned what those bands were, and where they led, and why he would not name them to a boy on his first morning. I will not spoil that for you either, Paul. It is enough to say that the refusal did more work than an answer would have. Beneath my frustration something took root that was not knowledge, not yet, only a disturbance continuing after the hand had gone.\par

\ornament

He reaches the balcony, the bands at the edge of the metal plate --- pale, dark, pale, dark --- and the old man who would not name them, and the old want moves in me before I can put my hand on it. Because I know what they are, and I know, better than he did at fifteen, where a boy who chases them arrives. I lean in. I do not mean to; the machine is between us, and I have leaned across it the width of a hand before I catch myself.\par

``You said it would be years before you understood those,'' I say. Level. Only curious. ``And where they led. But you did learn, in the end. Both.''\par

``I did.'' He watches me the way I taught his younger self to watch a lamp --- not the brightness, but the hand that placed it.\par

``Then we will come to it, when the account reaches it,'' I say --- and I hear how quickly I have said it, and so does he. ``Where did they lead.''\par

He is quiet a moment. When he answers, he gives me the Mountain and not the man. ``The account first, Paul.'' My own rule, my own words, turned in his mouth and set gently back on the table between us. ``All of it, in order, and the questions after. I will not spoil the shape of my own life to keep a servant's record straight.''\par

I sit back. I let the reels turn, and I put my face in order, and I say, ``As you like. In order.'' A servant would have been relieved. I am not certain I remembered to seem relieved.\par

He does not accuse me. He only lifts one finger from the table --- the small thing he does when he has set a matter aside to keep --- and turns the reels on. I taught him that a man is known by the light he reaches for. This morning I reached, and he saw the hand.\par

\ornament

At the candidate passage he stopped. ``You will come again tomorrow. Earlier.''\par

``Will I repeat the rooms?''\par

``No.''\par

``Then how do I prepare?''\par

``By noticing what you failed to notice today.''\par

``That is not preparation.''\par

``It is the only preparation Light respects.'' He turned to go, then paused. ``Sa-Kailan.'' The prefix still startled me. ``A Candidate learns where light falls.''\par

``A Veil learns who placed it,'' I said.\par

``And a Keeper?''\par

``Learns what must remain dark.''\par

He held my gaze. ``Do not rush to the third sentence because the first wounds your pride.''\par

Then he was gone down the eastern passage, his robe catching one pale shaft and losing it at the turn.\par

\ornament

I took the longer way back to my chamber. Twice I stopped to examine lamps I had passed a hundred times; once I crouched to find where a floor reflection began and found instead a servant watching me with polite concern. The Mountain had not changed. I had. That is the whole of what that first day did, and it was enough, and I want you to notice that nothing forbidden had happened yet. They had taught me to see. It is the last gift a cage should give a prisoner, and they gave it to me freely, and called it obedience.\par

In my chamber I unwrapped the thumb. The cut had reopened during the exercises, a dark crescent on the cloth. I sat beside the basin until the water stilled, then leaned over it and looked at a boy with a new prefix and an old hunger. I touched the surface with one finger, and my face broke into rings that met the basin's edge and returned inward, crossing one another into patterns I could neither keep nor understand. Beside the basin lay the stone token from the mirror chamber. I turned it once between finger and thumb and let it slip into the water.\par

It struck the side, turned, and settled near the bottom with the sign of Light facing up. I reached for it, and stopped.\par

The circle cut through the token no longer looked like itself. Its edge thickened when I moved my head left and narrowed when I moved right. The aperture seemed to slide away from the stone that held it. I had seen such things all my life, a spoon bent in a cup, a hand made strange in a washing pool, lamp-lines trembling through jars of oil. Everyone had seen them. Seeing was not the same as noticing.\par

I lowered my face nearer the surface. The lamp above the basin curved along the water's edge and broke against the rounded wall, and the light did not only return. It entered, turned, thickened, and came back altered. The basin was not a mirror, not exactly, nor a lens, though it borrowed something from lenses. Water did not merely remember disturbance. It changed the path of sight.\par

A brightness went through me that had nothing to do with the lamp. I looked to the door, half expecting Va-Raedin to be standing there already, knowing, already prepared to reduce discovery into correction. The chamber was empty. No Veil watched me. No Keeper named the lesson. No doctrine told me where to place my eye. Perhaps everyone in Light knew this. Perhaps he had known it for years and tomorrow would make me feel the smallness of it.\par

But perhaps not everyone had looked at water in this way.\par

I moved the token with two fingers and watched the sign slide one way while the stone moved another, and the excitement sharpened until it felt almost indecent. I forgot my hunger, my knees, my thumb, the Hall, Sheva, all of it, for a few breaths. Light entered water and did not remain obedient. That was the thing. That was the crack.\par

I withdrew my hand and waited for the surface to calm. Light did not simply reveal what was hidden. It could be delayed, diverted, bent by what it crossed. Water was not only a thing seen by light. Water could change what seeing meant. I did not yet know what that meant.\par

That made it better.\par

And that, Paul, if you want a beginning after all, is nearer to one than anything in the Tomb. Not the day they cut me. The night I found something no one had handed me, and looked at the door, and was glad no one had seen. A man who has learned to be glad the door is empty has already begun, in the smallest possible way, to carry something out of the room.\par

\ornament

He finishes on the boy at the basin --- the token turning in the water, the sign sliding off the stone, the empty doorway he was glad to find empty --- and for a while neither of us speaks. The reels turn. He has just told me, without knowing he has told me, the hour he became a man who could carry a thing out of a room. He offers it as the beginning I asked for. He does not hear it as a confession. I do.\par

Then he looks up, past me, at the lamp on the wall behind my chair. He has sat beneath it four days and not once asked. Now, softened by his own memory and off his guard, he asks.\par

``There is no flame in it,'' he says. ``I have watched. No oil, no wick, nothing to tend, and it does not gutter. What feeds it.''\par

``A kind of moving water we have learned to carry in wire,'' I say. True enough, and true is cheap, and the cheap answer keeps him easy. ``It runs, and the running makes the light. No one tends it. It does not know that it is watched.''\par

He considers the lamp a long moment --- the light with no keeper, the flame that answers to no hand in the room. ``A whole people who no longer guard their fire,'' he says at last, and I cannot tell whether it is wonder or pity in his voice, and I do not ask, because the difference is the one thing in this room that is his and not mine.\par

``He told you to come earlier the next day,'' I say. ``Into Light again. Begin there.''\par

He nods. He does not stop the tape. He has learned, this morning, to be glad of an empty doorway. He has not yet learned to be afraid of a full reel.\par

\ornament
